\chapter{Layout e Responsive Design}

Nel capitolo precedente abbiamo visto che cos'è il CSS, come includerlo, alcune
proprietà di base, i selettori, l'algoritmo della cascata e l'ereditarietà.
Adesso passiamo al problema più concreto: \textbf{disporre gli elementi nello
spazio}.

\section{Unità di misura}

Diverse proprietà CSS servono a stabilire la dimensione degli elementi o delle
loro aree: \attr{width}, \attr{height}, \attr{font-size}, \attr{margin},
\attr{padding}, \attr{border} e altre.

Nel dimensionare gli elementi si possono usare \term{lunghezze assolute} oppure
\term{lunghezze relative} (che dipendono da altre dimensioni, come quella del
carattere o della finestra).

\subsection{Unità assolute}

Le lunghezze assolute si esprimono con un numero seguito da un'unità di misura.

\begin{affianco}[0.55]
\begin{lstlisting}[style=css, name=assolute.css]
div {
  background: red;
  width: 2.54cm;             /* centimetri */
  height: 1in;               /* pollici */
  border: 2mm solid black;   /* millimetri */
  color: white;
  font-size: 24px;           /* pixel */
}
\end{lstlisting}
\begin{lstlisting}[style=html, name=assolute.html]
<div>hello</div>
\end{lstlisting}
\risultato
\begin{browser}[assolute.html]
  \begin{tikzpicture}
    \fill[css-red] (0,0) rectangle (2.54,2.54);
    \draw[black, line width=2mm] (-0.1,-0.1) rectangle (2.64,2.64);
    \node[anchor=north west, text=white, font=\Lora\fontsize{18}{20}\selectfont]
      at (0.05,2.49) {hello};
  \end{tikzpicture}
\end{browser}
\wtnota{Il riquadro è disegnato in scala reale: se stampi questa pagina,
        misura davvero 2,54\,cm di lato.}
\end{affianco}

\warning{Il pixel CSS non è il pixel dello schermo}{
  Un \code{px} in CSS non corrisponde necessariamente a un pixel fisico del
  monitor. Su uno schermo ad alta densità un pixel CSS può valere due o tre
  pixel fisici: è il motivo per cui un testo di 16px non diventa illeggibile su
  un telefono moderno. Le unità fisiche (\code{cm}, \code{mm}, \code{in}) hanno
  senso soltanto per la stampa; su schermo sono definite in rapporto fisso al
  pixel CSS e non misurano davvero un centimetro.}

\subsection{Percentuali}

Le percentuali definiscono le dimensioni in rapporto all'\textbf{elemento
padre}. Si possono usare con molte proprietà, fra cui \attr{width},
\attr{height}, \attr{font-size}, \attr{margin} e \attr{padding}.

\begin{esempio}{Dimensioni percentuali}
\begin{affianco}[0.5]
\begin{lstlisting}[style=css, name=percentuali.css]
.a {
  width: 300px; height: 200px;
  background: red;
}

.b {
  width: 50%;     /* 150px */
  height: 50%;    /* 100px */
  background: blue;
}
\end{lstlisting}

\begin{lstlisting}[style=html, name=percentuali.html]
<div class="a">
  <div class="b"></div>
</div>
\end{lstlisting}
\risultato
\begin{browser}[percentuali.html]
  \begin{tikzpicture}[font=\Lato\scriptsize,
      lbl/.style={text=white, anchor=south west, font=\Lato\scriptsize}]
    \fill[css-red]  (0,0)   rectangle (4.5,3);
    \fill[css-blue] (0,1.5) rectangle (2.25,3);
    \node[lbl] at (0.12,1.62) {\textbf{.b} \ 150$\times$100};
    \node[lbl] at (0.12,0.14) {\textbf{.a} \ 300$\times$200};
  \end{tikzpicture}
\end{browser}
\wtnota{Il figlio occupa metà della larghezza e metà dell'altezza del padre.}
\end{affianco}
\end{esempio}

\subsection{Unità relative al carattere}

% Nota: le celle della colonna p iniziano con \mbox{} perche' una cella che
% comincia direttamente con \code (un box in font monospaziato) sposta in
% basso la prima riga rispetto alla colonna accanto.
\begin{center}
\renewcommand{\arraystretch}{1.45}
\begin{tabular}{@{}l@{\hspace{1.1cm}}>{\raggedright\arraybackslash}p{9.6cm}@{}}
  \textbf{Unità} & \textbf{Significato} \\
  \code{em}  & \mbox{}\code{1em} è la dimensione del carattere \textbf{corrente};
               \code{1.5em} è il 50\% più grande. Storicamente, in tipografia,
               la dimensione del carattere derivava dalla larghezza della
               \guillemotleft M\guillemotright\ maiuscola: da qui il nome. \\
  \code{rem} & \mbox{}\code{1rem} è la dimensione del carattere calcolata
               sull'elemento \textbf{radice} (per impostazione predefinita
               16px). \\
\end{tabular}
\end{center}

\begin{esempio}{\code{em} contro \code{rem}}
\begin{affianco}[0.55]
\begin{lstlisting}[style=css, name=relative.css]
p {
  font-family: sans-serif;
  font-size: 1.5rem;    /* 1.5 x 16px = 24px */
}

span {
  font-size: 1.5em;     /* 1.5 x 24px = 36px */
}
\end{lstlisting}

\begin{lstlisting}[style=html, name=relative.html]
<p>HELLO <span>SIZES</span></p>
\end{lstlisting}
\risultato
\begin{browser}[relative.html]
  {\Lato\fontsize{18}{22}\selectfont HELLO
   {\fontsize{27}{30}\selectfont SIZES}}
\end{browser}
\wtnota{\code{1.5rem} = 1,5 $\times$ 16px = \textbf{24px}; dentro al paragrafo,
        \code{1.5em} = 1,5 $\times$ 24px = \textbf{36px}.}
\end{affianco}

La differenza è tutta qui: \code{rem} si riferisce sempre alla radice del
documento, quindi è \textbf{stabile}; \code{em} si riferisce al contesto, quindi
si \textbf{accumula}. Annidando tre elementi con \code{font-size: 1.5em} si
ottiene $1{,}5^3 \approx 3{,}4$ volte la dimensione di partenza, effetto quasi
sempre indesiderato. Per questo, di norma, si usa \code{rem} per le dimensioni
del testo e si riserva \code{em} ai casi in cui si vuole davvero che una
misura scali insieme al testo circostante (per esempio il \attr{padding} di un
pulsante).
\end{esempio}

\subsection{Unità relative al viewport}

Il \term{viewport} è l'area della finestra del browser in cui viene mostrato il
documento.

\begin{itemize}
  \item \code{vw}: \code{1vw} è l'1\% della \textbf{larghezza} del viewport;
  \item \code{vh}: \code{1vh} è l'1\% dell'\textbf{altezza} del viewport.
\end{itemize}

\begin{affianco}[0.55]
\begin{lstlisting}[style=css, name=viewport.css]
div { height: 50vh; }     /* meta' dell'altezza */

.a { background: red; }

.b {
  background: blue; color: white;
  width: 50vw;            /* meta' della larghezza */
  font-size: 10vw;        /* scala con la finestra! */
}
\end{lstlisting}
\begin{lstlisting}[style=html, name=viewport.html]
<div class="a"></div>
<div class="b">HELLO</div>
\end{lstlisting}
\risultato
\begin{browser}[viewport.html]
  \begin{tikzpicture}[font=\Lato\scriptsize]
    \fill[css-red]  (0,1.6) rectangle (6.6,3.2);
    \fill[css-blue] (0,0)   rectangle (3.3,1.6);
    \node[anchor=west, text=white, font=\Lora\fontsize{16}{18}\selectfont]
      at (0.1,0.9) {HELLO};
    \draw[{Stealth[length=4pt]}-{Stealth[length=4pt]}, draw=neutral!45!black,
          line width=0.5pt] (0,3.42) -- (6.6,3.42);
    \node[anchor=south, text=neutral!35!black] at (3.3,3.45) {100vw = larghezza finestra};
    \draw[{Stealth[length=4pt]}-{Stealth[length=4pt]}, draw=neutral!45!black,
          line width=0.5pt] (6.82,1.6) -- (6.82,3.2);
    \node[anchor=west, text=neutral!35!black, rotate=-90] at (6.9,3.2) {50vh};
  \end{tikzpicture}
\end{browser}
\wtnota{Allargando la finestra, tutto cresce con lei: riquadri \textit{e}
        dimensione del testo.}
\end{affianco}

\section{Il box model}

\dfn{Box model}{
  In CSS \textbf{ogni elemento HTML è una scatola} (\textit{box}). Ogni scatola
  è composta da quattro aree concentriche: \term{content box}, \term{padding
  box}, \term{border box} e \term{margin box}.}

\begin{figure}[H]
  \centering
  \begin{tikzpicture}[font=\Lato\small,
      area/.style={font=\Lato\small\bfseries},
      mis/.style={{Stealth[length=4pt]}-{Stealth[length=4pt]},
                  draw=neutral!45!black, line width=0.6pt},
      mist/.style={font=\Lato\scriptsize, text=neutral!30!black, fill=white,
                   inner sep=1.5pt}]

    % margin box
    \fill[orange-gradient-1!40!white] (0,0) rectangle (13,7.2);
    \draw[dashed, draw=orange-gradient-6, line width=0.9pt] (0,0) rectangle (13,7.2);
    \node[area, anchor=north west, text=orange-gradient-6!65!black]
      at (0.18,7.02) {margin box};

    % border box
    \fill[neutral!30!white] (1.5,1.1) rectangle (11.5,6.1);
    \draw[draw=neutral!45!black, line width=1.6pt] (1.5,1.1) rectangle (11.5,6.1);
    \node[area, anchor=north east, text=neutral!20!black]
      at (11.32,5.92) {border box};

    % padding box
    \fill[green-gradient-1!55!white] (2.5,1.9) rectangle (10.5,5.3);
    \draw[dashed, draw=green-gradient-6, line width=0.8pt] (2.5,1.9) rectangle (10.5,5.3);
    \node[area, anchor=south east, text=green-gradient-6!55!black]
      at (10.32,2.05) {padding box};

    % content box
    \fill[primary!16!white] (3.9,2.7) rectangle (9.1,4.5);
    \draw[draw=primary!65!black, line width=0.9pt] (3.9,2.7) rectangle (9.1,4.5);
    \node[area, text=primary!70!black] at (6.5,3.6) {content box};

    % misure, ognuna su una riga diversa per non sovrapporsi
    % le tre misure, a scaletta, per non sovrapporsi fra loro
    \draw[mis] (0.75,0) -- (0.75,1.1);
    \node[mist, anchor=east] at (0.65,0.55) {margin};
    \draw[mis] (2.0,1.1) -- (2.0,1.9);
    \node[mist, anchor=west] at (2.12,1.5) {border};
    \draw[mis] (3.2,1.9) -- (3.2,2.7);
    \node[mist, anchor=west] at (3.32,2.3) {padding};
  \end{tikzpicture}
  \caption{Le quattro aree del box model.}
  \label{fig:box-model}
\end{figure}

I valori di ciascuna area si leggono comodamente nei DevTools: selezionando un
elemento, il pannello mostra il box model con le misure effettive, area per
area.

\begin{figure}[H]
  \centering
  \begin{tikzpicture}[font=\FiraCodeNL\scriptsize,
      etic/.style={font=\Lato\scriptsize\bfseries, anchor=north west},
      num/.style={font=\FiraCodeNL\scriptsize, text=neutral!20!black}]

    \fill[orange-gradient-1!45!white] (0,0) rectangle (9,5.2);
    \draw[draw=orange-gradient-6, line width=0.7pt] (0,0) rectangle (9,5.2);
    \node[etic, text=orange-gradient-6!65!black] at (0.12,5.08) {margin};

    \fill[yellow-gradient-1!55!white] (0.95,0.75) rectangle (8.05,4.45);
    \draw[draw=yellow-gradient-6, line width=0.7pt] (0.95,0.75) rectangle (8.05,4.45);
    \node[etic, text=yellow-gradient-6!60!black] at (1.07,4.33) {border};

    \fill[green-gradient-1!55!white] (1.85,1.45) rectangle (7.15,3.75);
    \draw[draw=green-gradient-6, line width=0.7pt] (1.85,1.45) rectangle (7.15,3.75);
    \node[etic, text=green-gradient-6!55!black] at (1.97,3.63) {padding};

    \fill[primary!16!white] (2.9,2.1) rectangle (6.1,3.1);
    \draw[draw=primary!60!black, line width=0.7pt] (2.9,2.1) rectangle (6.1,3.1);
    \node[num, text=primary!70!black] at (4.5,2.6) {100 $\times$ 92.8};

    % valori, uno per lato
    \node[num] at (4.5,4.85) {10};  \node[num] at (4.5,0.35) {10};
    \node[num] at (0.48,2.6) {10};  \node[num] at (8.52,2.6) {10};
    \node[num] at (4.5,4.1)  {20};  \node[num] at (4.5,1.1)  {20};
    \node[num] at (1.4,2.6)  {20};  \node[num] at (7.6,2.6)  {20};
    \node[num] at (4.5,3.43) {20};  \node[num] at (4.5,1.78) {20};
    \node[num] at (2.37,2.6) {20};  \node[num] at (6.63,2.6) {20};
  \end{tikzpicture}
  \caption{Lo stesso box model come lo mostrano i DevTools, con le misure in
           pixel: contenuto di 100$\times$92,8, poi 20 di padding, 20 di bordo e
           10 di margine su ogni lato.}
  \label{fig:box-model-devtools}
\end{figure}

\begin{itemize}
  \item il \term{content box} è dove vivono i figli dell'elemento;
  \item il \term{padding} (spaziatura \textit{interna}) separa il contenuto dal
        bordo;
  \item il \term{border} è il confine dell'elemento;
  \item il \term{margin} crea spazio \textit{attorno} all'elemento.
\end{itemize}

La dimensione di ciascuna area si controlla con le dichiarazioni CSS
corrispondenti. Il comportamento e l'aspetto delle scatole dipendono da tre
fattori: la \textbf{modalità di layout}, il \textbf{contenuto} e le
\textbf{proprietà del box model}.

\info{L'overflow}{
  Il contenuto potrebbe essere troppo grande per stare nella sua scatola. La
  parte che non ci sta si chiama \term{overflow}, e la proprietà
  \attr{overflow} permette di decidere che cosa farne: mostrarla comunque
  (\code{visible}, il valore predefinito), tagliarla (\code{hidden}) o
  aggiungere delle barre di scorrimento (\code{scroll}, \code{auto}).}

\error{La sorpresa di \attr{width}}{
  Per impostazione predefinita, \attr{width} e \attr{height} dimensionano il
  \textbf{solo content box}. Un elemento con \code{width: 200px},
  \code{padding: 20px} e \code{border: 5px} occupa quindi
  $200 + 20 \cdot 2 + 5 \cdot 2 = 250$ pixel: non 200. È il singolo aspetto
  del CSS che genera più confusione. La soluzione universalmente adottata è
  dichiarare in cima al foglio di stile

  \code{* \{ box-sizing: border-box; \}}

  che fa sì che \attr{width} indichi la larghezza \textbf{fino al bordo}
  incluso. Da quel momento 200px significa davvero 200px.}

\section{Il flow layout}

Per impostazione predefinita, le scatole vengono disposte secondo il \term{flow
layout} (o \textit{normal flow}).

\begin{itemize}
  \item gli elementi \term{inline} si dispongono uno accanto all'altro nella
        \textbf{direzione inline}, come le parole di una frase;
  \item gli elementi \term{block} si dispongono uno \textbf{sotto} l'altro.
\end{itemize}

\begin{figure}[H]
  \centering
  \begin{tikzpicture}[font=\Lato\scriptsize,
      il/.style={draw=primary!65!black, fill=primary!12!white, line width=0.7pt,
                 minimum height=0.55cm, inner xsep=6pt},
      bl/.style={draw=accent!60!black, fill=accent!12!white, line width=0.7pt,
                 minimum height=0.6cm, minimum width=7.4cm}]

    \node[il, anchor=north west] (i1) at (0,0)   {inline};
    \node[il, anchor=north west] (i2) at (1.6,0) {inline};
    \node[il, anchor=north west] (i3) at (0,-0.75)   {inline};
    \node[il, anchor=north west] (i4) at (1.6,-0.75) {inline};
    \node[bl, anchor=north west] (b1) at (0,-1.55) {block};
    \node[il, anchor=north west] (i5) at (0,-2.35) {inline};

    \draw[-{Stealth[length=5pt]}, draw=primary!65!black, line width=0.8pt]
      (0,0.35) -- (3.1,0.35);
    \node[font=\Lato\scriptsize, text=primary!65!black, anchor=south west]
      at (0,0.42) {direzione inline};

    \draw[-{Stealth[length=5pt]}, draw=accent!60!black, line width=0.8pt]
      (7.75,0) -- (7.75,-2.9);
    \node[font=\Lato\scriptsize, text=accent!60!black, anchor=west, rotate=-90]
      at (7.95,-0.1) {direzione block};
  \end{tikzpicture}
  \caption{Il flusso normale: gli inline in riga, i block uno sotto l'altro.}
  \label{fig:flow}
\end{figure}

\subsection{La proprietà display}

\begin{center}
\renewcommand{\arraystretch}{1.45}
\begin{tabular}{@{}l@{\hspace{1.0cm}}>{\raggedright\arraybackslash}p{9.4cm}@{}}
  \textbf{Valore} & \textbf{Effetto} \\
  \code{inline} & posiziona la scatola accanto alla precedente, nella direzione
    inline. \attr{width} e \attr{height} vengono \textbf{ignorate}. È il valore
    predefinito di \tg{a}, \tg{em} e \tg{strong}. \\
  \code{block} & posiziona la scatola su una riga dedicata. Se non specificato
    diversamente, i block si espandono per tutta la dimensione inline
    disponibile. È il valore predefinito di \tg{p} e delle intestazioni. \\
  \code{inline-block} & come \code{inline}, ma consente di impostare
    \attr{width} e \attr{height}. \\
  \code{none} & rimuove completamente l'elemento dalla visualizzazione. \\
\end{tabular}
\end{center}

% Scatola gialla con bordo, come l'em degli esempi su display
\newcommand{\demq}[1]{\fcolorbox{black}{css-yellow}{%
  \begin{minipage}[c][1.15cm][t]{1.15cm}\raggedright\Lora\small #1\end{minipage}}}
\newcommand{\demi}[1]{{\setlength{\fboxsep}{1pt}%
  \fcolorbox{black}{css-yellow}{\Lora\small\textit{#1}}}}

\begin{esempio}{Lo stesso markup con quattro display diversi}
\begin{affianco}[0.52]
\begin{lstlisting}[style=html, name=display.html]
<h3>Intestazione</h3>
<p>
  Sto imparando <em>HTML</em> e <em>CSS</em>
</p>
\end{lstlisting}
\begin{lstlisting}[style=css, name=display.css]
em {
  display: inline;      /* poi: block,
                           inline-block, none */
  background: yellow;
  border: 1px solid;
  width: 50px; height: 50px;
}
\end{lstlisting}
\risultato
\begin{browser}[display: inline]
  {\Lora\bfseries\large Intestazione}\\[5pt]
  Sto imparando \demi{HTML} e \demi{CSS}
\end{browser}
\begin{browser}[display: block]
  {\Lora\bfseries\large Intestazione}\\[5pt]
  Sto imparando\\[2pt]
  \demq{\textit{HTML}}\\[2pt]
  e\\[2pt]
  \demq{\textit{CSS}}
\end{browser}
\begin{browser}[display: inline-block]
  {\Lora\bfseries\large Intestazione}\\[5pt]
  Sto imparando \demq{\textit{HTML}} e \demq{\textit{CSS}}
\end{browser}
\begin{browser}[display: none]
  {\Lora\bfseries\large Intestazione}\\[5pt]
  Sto imparando e
\end{browser}
\end{affianco}

\begin{center}
\renewcommand{\arraystretch}{1.45}
\begin{tabular}{@{}l@{\hspace{1.0cm}}>{\raggedright\arraybackslash}p{12.8cm}@{}}
  \textbf{Valore} & \textbf{Che cosa si vede} \\
  \code{inline} & i due \tg{em} restano nella frase; \attr{width} e
    \attr{height} sono \textbf{ignorate}, quindi le scatole gialle sono larghe
    quanto le parole. \\
  \code{block} & ogni \tg{em} va su una riga propria e diventa un quadrato di
    50$\times$50, spezzando la frase in tre pezzi. \\
  \code{inline-block} & i due \tg{em} restano nella frase \textbf{e} diventano
    quadrati di 50$\times$50. \\
  \code{none} & i due \tg{em} spariscono: si legge soltanto
    \guillemotleft Sto imparando e\guillemotright, con lo spazio che si richiude. \\
\end{tabular}
\end{center}
\end{esempio}

\warning{\code{display: none} contro \code{visibility: hidden}}{
  \code{display: none} toglie l'elemento dal flusso: non occupa più spazio, come
  se non esistesse. \code{visibility: hidden} lo rende invisibile ma
  \textbf{conserva} lo spazio che occupava. Nessuno dei due, però, rimuove
  l'elemento dal documento: il contenuto viaggia comunque sulla rete e resta
  leggibile nel sorgente.}

\section{Float e clear}

La proprietà \attr{float} fa \guillemotleft galleggiare\guillemotright\ un elemento nella direzione
indicata, e fa sì che gli elementi fratelli successivi gli si dispongano
attorno.

\begin{esempio}{Un'immagine con il testo che le scorre attorno}
\begin{lstlisting}[style=html, name=float.html]
<h1>Fu-Tzu</h1>
<img src="pic.jpg">
<p>Con una lunga barba e vesti ornate di argute battute sul codice,
Fu-Tzu era una leggenda fra i programmatori.</p>
<p>Si dice che la sua prima parola sia stata <code>printf</code>. Parte
della sua saggezza fu trascritta nel
<a href="index.html">Libro della Programmazione</a>.</p>
\end{lstlisting}

\begin{affianco}[0.42]
\begin{lstlisting}[style=css, name=float.css]
img {
  width: 20vw;
  min-width: 100px;
  float: left;
  margin-right: 1rem;
}
\end{lstlisting}
\risultato
\begin{browser}[float.html]
  {\Lora\bfseries\Large Fu-Tzu}\\[4pt]
  \begin{minipage}[t]{2.3cm}\vspace{0pt}
    \includegraphics[width=2.3cm]{sources/fu-tzu.jpg}
  \end{minipage}\hspace{0.28cm}%
  \begin{minipage}[t]{6.0cm}\vspace{0pt}\raggedright\Lora\small
    Con una lunga barba e vesti ornate di argute battute sul codice, Fu-Tzu era
    una leggenda fra i programmatori.\par\vspace{4pt}
    Si dice che la sua prima parola sia stata \texttt{\small printf}. Parte
  \end{minipage}\par\vspace{2pt}
  {\raggedright\Lora\small della sua saggezza fu trascritta nel
   \wtlink{Libro della Programmazione}.\par}
\end{browser}
\wtnota{Il testo scorre a fianco dell'immagine e, quando questa finisce,
        riprende a tutta larghezza.}
\end{affianco}

L'immagine si sposta a sinistra e i due paragrafi le scorrono a fianco.

La proprietà \attr{clear} serve a impedire che un fratello successivo si
disponga accanto all'elemento flottante:

\begin{affianco}[0.42]
\begin{lstlisting}[style=css]
p + p { clear: both; }
\end{lstlisting}
\risultato
\begin{browser}[float.html con clear]
  {\Lora\bfseries\Large Fu-Tzu}\\[4pt]
  \begin{minipage}[t]{2.3cm}\vspace{0pt}
    \includegraphics[width=2.3cm]{sources/fu-tzu.jpg}
  \end{minipage}\hspace{0.28cm}%
  \begin{minipage}[t]{6.0cm}\vspace{0pt}\raggedright\Lora\small
    Con una lunga barba e vesti ornate di argute battute sul codice, Fu-Tzu era
    una leggenda fra i programmatori.
  \end{minipage}\par\vspace{4pt}
  {\raggedright\Lora\small Si dice che la sua prima parola sia stata
   \texttt{\small printf}. Parte della sua saggezza fu trascritta nel
   \wtlink{Libro della Programmazione}.\par}
\end{browser}
\wtnota{Il secondo paragrafo aspetta la fine dell'immagine e comincia sotto.}
\end{affianco}

Con questa regola il secondo paragrafo comincia \textbf{sotto} l'immagine
anziché al suo fianco.
\end{esempio}

\info{Il float oggi}{
  \attr{float} nacque per un compito tipografico preciso: far scorrere il testo
  attorno a un'immagine. Per anni, in mancanza di alternative, è stato però
  piegato a costruire interi impianti di pagina a colonne, con una quantità di
  trucchi (i famigerati \textit{clearfix}) per rimediare ai suoi effetti
  collaterali. Con Flexbox e Grid quell'uso non ha più ragione di esistere:
  \attr{float} è tornato a fare ciò per cui era nato.}

\section{Posizionamento}

La proprietà \attr{position} cambia il modo in cui un elemento si comporta nel
flusso normale del documento. Il valore predefinito è \code{static}.

\begin{center}
\renewcommand{\arraystretch}{1.45}
\begin{tabular}{@{}l@{\hspace{1.0cm}}>{\raggedright\arraybackslash}p{9.4cm}@{}}
  \textbf{Valore} & \textbf{Comportamento} \\
  \code{static}   & flusso normale (predefinito). \\
  \code{relative} & posizionato rispetto alla \textbf{propria posizione
                    normale}. \\
  \code{absolute} & posizionato rispetto all'\textbf{antenato posizionato} più
                    vicino. \\
  \code{fixed}    & posizionato rispetto al \textbf{viewport}. \\
  \code{sticky}   & posizionato in base alla \textbf{posizione di scorrimento}
                    dell'utente. \\
\end{tabular}
\end{center}

\subsection{relative}

\begin{affianco}[0.5]
\begin{lstlisting}[style=css, name=relative.css]
div {
  position: relative;
  top: 50px; left: 50px;
  background: red;
  width: 100px; height: 100px;
}

div div {
  background: blue;
}
\end{lstlisting}
\begin{lstlisting}[style=html, name=relative.html]
<div><div></div></div>
\end{lstlisting}
\risultato
\begin{browser}[relative.html]
  \begin{tikzpicture}[font=\Lato\scriptsize]
    \draw[dashed, draw=neutral!50!black, line width=0.6pt] (0,1.5) rectangle (1.5,3);
    \node[text=neutral!40!black, anchor=south west, align=left]
      at (0.05,3.1) {posizione normale\\(spazio riservato)};
    \fill[css-red]  (0.75,0.75) rectangle (2.25,2.25);
    \fill[css-blue] (0.75,0.75) rectangle (2.25,2.25);
    \draw[{Stealth[length=4pt]}-, draw=accent!70!black, line width=0.7pt]
      (0.78,2.22) -- (0.05,2.95);
    \node[text=accent!70!black, anchor=west] at (2.4,1.5) {spostato di 50px};
  \end{tikzpicture}
\end{browser}
\wtnota{Il figlio, che non ha dimensioni proprie, riempie il padre: si vede
        solo il blu.}
\end{affianco}

L'elemento viene spostato di 50 pixel verso il basso e verso destra rispetto a
dove sarebbe finito normalmente. Lo spazio che occupava nel flusso, però,
\textbf{resta riservato}: gli altri elementi non si spostano per riempirlo.

\subsection{absolute}

Gli elementi con \code{position: absolute} sono posizionati rispetto
all'antenato posizionato (cioè non \code{static}) più vicino. Se non ne esiste
alcuno, il riferimento è il documento.

A differenza di \code{relative}, un elemento assoluto viene \textbf{tolto dal
flusso}: lo spazio che occupava viene liberato e gli altri elementi si
ridispongono come se non ci fosse.

\begin{affianco}[0.45]
\begin{lstlisting}[style=css, name=absolute.css]
.a { position: relative; }

.c {
  position: absolute;
  bottom: 10px; right: 10px;
}
\end{lstlisting}
\begin{lstlisting}[style=html, name=absolute.html]
<div class="a">
  A: questo div ha position: relative
  <div class="b">
    B: questo div non ha position
    <div class="c">
      C: position absolute
    </div>
  </div>
</div>
\end{lstlisting}
\risultato
\begin{browser}[absolute.html]
  \begin{tikzpicture}[font=\Lora\scriptsize]
    \fill[accent!12!white] (0,0) rectangle (8,3.4);
    \node[anchor=north west, align=left] at (0.1,3.3) {A: questo div ha \code{position: relative}};
    \fill[azure-gradient-1!45!white] (0.35,0.45) rectangle (7.65,2.75);
    \node[anchor=north west] at (0.45,2.65) {B: questo div non ha \code{position}};
    \fill[red-gradient-1!45!white] (4.1,0.18) rectangle (7.82,0.92);
    \node[anchor=west] at (4.2,0.55) {C: \code{position: absolute}};
    \draw[{Stealth[length=4pt]}-{Stealth[length=4pt]}, draw=neutral!45!black,
          line width=0.5pt] (7.82,0.55) -- (8,0.55);
    \node[font=\Lato\tiny, text=neutral!35!black, anchor=south] at (7.9,0.6) {10px};
  \end{tikzpicture}
\end{browser}
\wtnota{\code{.c} si posiziona rispetto ad \code{.a}, l'antenato posizionato
        più vicino, non rispetto a \code{.b}.}
\end{affianco}

\subsection{fixed}

\begin{affianco}[0.45]
\begin{lstlisting}[style=css, name=fixed.css]
a {
  position: fixed;
  bottom: 2rem; right: 2rem;
  background-color: white;
  border: 2px solid;
  padding: 1em;
}
\end{lstlisting}
\risultato
\begin{browser}[fixed.html]
  \begin{tikzpicture}[font=\Lato\scriptsize]
    \node[anchor=north west, font=\Lora\bfseries\large] at (0,3.3) {Saggezza};
    \foreach \y in {2.7,2.45,2.2,1.95,1.7,1.45,1.2,0.95,0.7,0.45}
      \fill[neutral!22!white] (0,\y) rectangle (7.2,\y+0.13);
    \draw[draw=black, fill=white, line width=1pt] (5.2,0.25) rectangle (7.0,0.85);
    \node[font=\Lora\scriptsize] at (6.1,0.55) {Torna su};
    \draw[-{Stealth[length=5pt]}, draw=accent!70!black, line width=0.8pt]
      (7.6,2.6) -- (7.6,0.6);
    \node[text=accent!70!black, anchor=west, rotate=-90, align=center]
      at (7.75,2.6) {la pagina scorre,\\il pulsante no};
  \end{tikzpicture}
\end{browser}
\end{affianco}

L'elemento è posizionato rispetto al viewport e \textbf{non si muove} con lo
scorrimento della pagina. È il meccanismo dietro ai pulsanti \guillemotleft torna
su\guillemotright, alle barre di navigazione sempre visibili e ai banner dei cookie.

\subsection{sticky}

Gli elementi con \code{position: sticky} sono posizionati in base alla posizione
di scorrimento: sono una sorta di ibrido fra \code{relative} e \code{fixed}.

\begin{itemize}
  \item si comportano come \code{relative} finché non superano una soglia;
  \item da quel momento si comportano come \code{fixed}, fino a raggiungere il
        confine del proprio elemento padre.
\end{itemize}

\begin{affianco}[0.45]
\begin{lstlisting}[style=css, name=sticky.css]
nav {
  padding: 1em;
  background: gainsboro;
  position: sticky;
  top: 0px;    /* la soglia: si "incolla"
                  al bordo superiore */
}
\end{lstlisting}
\risultato
\begin{browser}[prima di scorrere]
  \begin{tikzpicture}[font=\Lato\scriptsize]
    \node[anchor=north west, font=\Lora\bfseries] at (0,1.5) {Il libro della programmazione};
    \fill[neutral!22!white] (0,0.95) rectangle (7.2,1.25);
    \node[anchor=west, font=\Lora\scriptsize, text=css-link] at (0.1,1.1)
      {Home \quad Saggezza \quad Dati e Controllo};
    \foreach \y in {0.6,0.35,0.1} \fill[neutral!18!white] (0,\y) rectangle (7.2,\y+0.13);
  \end{tikzpicture}
\end{browser}
\begin{browser}[dopo lo scorrimento]
  \begin{tikzpicture}[font=\Lato\scriptsize]
    \fill[neutral!22!white] (0,1.2) rectangle (7.2,1.5);
    \node[anchor=west, font=\Lora\scriptsize, text=css-link] at (0.1,1.35)
      {Home \quad Saggezza \quad Dati e Controllo};
    \foreach \y in {0.85,0.6,0.35,0.1} \fill[neutral!18!white] (0,\y) rectangle (7.2,\y+0.13);
    \draw[-{Stealth[length=4pt]}, draw=accent!70!black, line width=0.7pt]
      (7.45,0.3) -- (7.45,1.3);
    \node[text=accent!70!black, anchor=west, rotate=-90] at (7.6,1.35) {resta in cima};
  \end{tikzpicture}
\end{browser}
\end{affianco}

È il modo standard per realizzare una barra di navigazione che scorre con la
pagina finché non arriva in cima, e lì rimane.

\section{Flexbox}

Oltre al flusso normale, il CSS moderno offre due ulteriori meccanismi di
layout: \term{Flexbox} e \term{Grid}.

\begin{itemize}
  \item \textbf{Flexbox} è pensato per layout \textbf{monodimensionali}: una
        riga oppure una colonna.
  \item \textbf{Grid} è pensato per layout \textbf{bidimensionali}: righe e
        colonne insieme.
\end{itemize}

\subsection{Contenitori flex}

Un contenitore flex si dichiara con \code{display: flex}. È un elemento di
livello blocco, i cui figli diventano \term{flex item}.

Ogni contenitore flex ha un \term{asse principale} (\textit{main axis}) e un
\term{asse trasversale} (\textit{cross axis}). L'asse principale si stabilisce
con la proprietà \attr{flex-direction} (il valore predefinito è \code{row});
l'asse trasversale è quello ortogonale.

\begin{figure}[H]
  \centering
  \begin{tikzpicture}[font=\Lato\scriptsize,
      cont/.style={draw=primary!65!black, fill=primary!5!white, line width=0.9pt},
      it/.style={draw=accent!60!black, fill=accent!14!white, line width=0.7pt}]

    % row
    \draw[cont] (0,0) rectangle (5.4,2.4);
    \foreach \x in {0.35,2.0,3.65} { \draw[it] (\x,0.7) rectangle (\x+1.4,1.7); }
    \node[font=\Lato\scriptsize\bfseries, text=primary!70!black, anchor=south]
      at (2.7,2.55) {\code{flex-direction: row}};
    \draw[-{Stealth[length=5pt]}, draw=primary!70!black, line width=0.9pt]
      (0.35,0.35) -- (5.05,0.35);
    \node[text=primary!70!black, anchor=north] at (2.7,0.28) {asse principale};
    \draw[-{Stealth[length=5pt]}, draw=accent!70!black, line width=0.9pt]
      (5.75,2.05) -- (5.75,0.35);
    \node[text=accent!70!black, anchor=west, rotate=-90] at (5.95,2.0) {asse trasversale};

    % column
    \begin{scope}[xshift=8.2cm]
      \draw[cont] (0,0) rectangle (2.6,2.4);
      \foreach \y in {0.3,1.0,1.7} { \draw[it] (0.55,\y) rectangle (2.05,\y+0.5); }
      \node[font=\Lato\scriptsize\bfseries, text=primary!70!black, anchor=south]
        at (1.3,2.55) {\code{flex-direction: column}};
      \draw[-{Stealth[length=5pt]}, draw=primary!70!black, line width=0.9pt]
        (0.28,2.15) -- (0.28,0.25);
      \node[text=primary!70!black, anchor=west, rotate=-90] at (0.1,2.1) {asse principale};
      \draw[-{Stealth[length=5pt]}, draw=accent!70!black, line width=0.9pt]
        (0.55,-0.3) -- (2.05,-0.3);
      \node[text=accent!70!black, anchor=north] at (1.3,-0.37) {asse trasversale};
    \end{scope}
  \end{tikzpicture}
  \caption{I due assi di un contenitore flex.}
  \label{fig:flex-axes}
\end{figure}

\subsection{Flex item}

I figli del contenitore flex diventano flex item, vengono disposti lungo l'asse
principale e supportano alcune proprietà specifiche.

\begin{center}
\renewcommand{\arraystretch}{1.45}
\begin{tabular}{@{}l@{\hspace{1.0cm}}>{\raggedright\arraybackslash}p{9.2cm}@{}}
  \textbf{Proprietà} & \textbf{Effetto} \\
  \attr{flex-grow}   & fa espandere l'elemento occupando tutto lo spazio
                       disponibile lungo l'asse principale. \\
  \attr{flex-shrink} & controlla se l'elemento può ridurre la propria dimensione
                       di base per stare nel contenitore. \\
  \attr{flex-wrap}   & si imposta sul \textbf{contenitore} e controlla che cosa
                       accade in caso di traboccamento: gli elementi possono
                       andare a capo su più righe. \\
\end{tabular}
\end{center}

\subsection{justify-content}

La proprietà \attr{justify-content} del contenitore specifica come gestire lo
spazio libero lungo l'\textbf{asse principale}.

\begin{figure}[H]
  \centering
  \begin{tikzpicture}[font=\Lato\scriptsize,
      cont/.style={draw=primary!65!black, fill=primary!4!white, line width=0.8pt},
      it/.style={draw=accent!60!black, fill=accent!16!white, line width=0.6pt}]

    \def\W{7.6}
    \foreach \lbl/\i/\pos in {%
        flex-start/0/{0.00,0.95,1.90},
        flex-end/1/{4.80,5.75,6.70},
        center/2/{2.42,3.37,4.32},
        space-between/3/{0.00,3.35,6.70},
        space-around/4/{0.82,3.35,5.88},
        space-evenly/5/{1.22,3.35,5.48}} {
      \begin{scope}[yshift=-\i*0.95cm]
        \node[font=\Lato\scriptsize, anchor=east, text=neutral!25!black]
          at (-0.3,0.35) {\code{\lbl}};
        \draw[cont] (0,0) rectangle (\W,0.7);
        \foreach \x in \pos { \draw[it] (\x,0.1) rectangle (\x+0.9,0.6); }
      \end{scope}
    }
  \end{tikzpicture}
  \caption{I valori più usati di \attr{justify-content}.}
  \label{fig:justify}
\end{figure}

\subsection{align-content e align-items}

Analogamente, \attr{align-content} specifica come gestire lo spazio libero lungo
l'\textbf{asse trasversale}; i valori più usati sono \code{flex-start},
\code{flex-end} e \code{center}. La proprietà \attr{align-items} allinea invece
i singoli elementi lungo lo stesso asse.

\begin{esempio}{Un contenitore flex completo}
\begin{affianco}[0.52]
\begin{lstlisting}[style=html, name=flex.html]
<div class="container">
  <div class="item special"><em>1</em></div>
  <div class="item"><em>2</em></div>
  <div class="item"><em>3</em></div>
</div>
\end{lstlisting}

\begin{lstlisting}[style=css, name=flex.css]
.container {
  width: 300px;
  height: 50px;
  background: rgb(255, 182, 182);
  display: flex;
  justify-content: space-around;
  align-items: center;
}

.item {
  height: 20px; width: 50px;
  padding: 5px;
  background: rgb(169, 169, 255);
}
\end{lstlisting}

\risultato
\begin{browser}[flex.html]
  \begin{tikzpicture}[font=\Lora\scriptsize]
    \fill[red!25!white] (0,0) rectangle (7.5,1.25);
    \foreach \x/\n in {0.65/1, 3.0/2, 5.35/3} {
      \fill[blue!30!white] (\x,0.32) rectangle (\x+1.5,0.93);
      \node[text=neutral!20!black] at (\x+0.75,0.62) {\textit{\n}};
    }
    \draw[{Stealth[length=3pt]}-{Stealth[length=3pt]}, draw=neutral!45!black,
          line width=0.4pt] (0,1.45) -- (7.5,1.45);
    \node[font=\Lato\tiny, text=neutral!35!black, anchor=south] at (3.75,1.48) {300px};
  \end{tikzpicture}
\end{browser}
\wtnota{Lo spazio libero è distribuito attorno ai tre elementi, che restano
        centrati in verticale nei 50px di altezza.}
\end{affianco}

Gli elementi sono distribuiti orizzontalmente con spazio uguale attorno a
ciascuno (\code{space-around}) e centrati verticalmente
(\code{align-items: center}).
\end{esempio}

\info{Uno strumento per esercitarsi}{
  Il sito \urlx{https://flexbox.help/} permette di visualizzare l'effetto delle
  proprietà flex in tempo reale ed è molto utile per prenderci confidenza.}

\section{Grid}

\term{Grid} è un layout \textbf{bidimensionale}, basato su righe e colonne.

\begin{itemize}
  \item i contenitori grid si dichiarano con \code{display: grid};
  \item i figli diretti dei contenitori grid diventano \term{grid item};
  \item il contenitore definisce \textbf{numero e dimensione} delle proprie
        righe e colonne.
\end{itemize}

\begin{esempio}{Una griglia a due colonne e tre righe}
\begin{affianco}[0.55]
\begin{lstlisting}[style=css, name=grid.css]
.container {
  display: grid;
  grid-template-columns: 1fr 100px;        /* due colonne */
  grid-template-rows: 20px 50px 100px;     /* tre righe */
  gap: 10px;
}
\end{lstlisting}

\begin{lstlisting}[style=html, name=grid.html]
<div class="container">
  <div>A</div><div>B</div><div>C</div>
  <div>D</div><div>E</div><div>F</div>
</div>
\end{lstlisting}

\risultato
\begin{browser}[grid.html]
  \begin{tikzpicture}[font=\Lora\scriptsize,
      c/.style={draw=primary!60!black, fill=primary!12!white, line width=0.7pt}]
    \foreach \r/\h/\y in {0/0.35/2.55, 1/0.7/1.65, 2/1.2/0.25} {
      \draw[c] (0,\y) rectangle (4.0,\y+\h);
      \draw[c] (4.3,\y) rectangle (5.9,\y+\h);
    }
    \node at (2.0,2.72)  {A}; \node at (5.1,2.72)  {B};
    \node at (2.0,2.0)   {C}; \node at (5.1,2.0)   {D};
    \node at (2.0,0.85)  {E}; \node at (5.1,0.85)  {F};
    \node[font=\Lato\scriptsize, text=neutral!35!black, anchor=north] at (2.0,-0.05) {\code{1fr}};
    \node[font=\Lato\scriptsize, text=neutral!35!black, anchor=north] at (5.1,-0.05) {\code{100px}};
  \end{tikzpicture}
\end{browser}
\wtnota{Le righe sono alte 20, 50 e 100 pixel; il \code{gap} lascia 10 pixel
        fra una cella e l'altra.}
\end{affianco}

Gli elementi vengono collocati \textbf{riga per riga}: A e B nella prima, C e D
nella seconda, E ed F nella terza.
\end{esempio}

\subsection{L'unità fr}

L'unità \code{fr} (\textit{fraction}) è un'unità relativa specializzata che
funziona \textbf{solo} nei layout grid. Rappresenta una lunghezza flessibile
corrispondente a una \textbf{quota dello spazio disponibile}, e si comporta in
modo simile alla proprietà \attr{flex}.

Per esempio, \code{grid-template-columns: 1fr 1fr 1fr;} definisce tre colonne
che si spartiscono in parti uguali lo spazio disponibile. Nell'esempio
precedente, invece, \code{1fr 100px} significa: la seconda colonna è larga
esattamente 100 pixel, la prima prende \textbf{tutto il resto}.

\subsection{Direzione di riempimento}

Per impostazione predefinita i grid item vengono collocati lungo le righe, ma è
possibile collocarli lungo le colonne con la proprietà
\code{grid-auto-flow: column}:

\begin{affianco}[0.55]
\begin{lstlisting}[style=css]
.container {
  display: grid;
  grid-template-columns: 1fr 100px;
  grid-template-rows: 20px 50px 100px;
  grid-auto-flow: column;
  gap: 10px;
}
\end{lstlisting}
\risultato
\begin{browser}[grid.html]
  \begin{tikzpicture}[font=\Lora\scriptsize,
      c/.style={draw=primary!60!black, fill=primary!12!white, line width=0.7pt}]
    \foreach \r/\h/\y in {0/0.35/2.55, 1/0.7/1.65, 2/1.2/0.25} {
      \draw[c] (0,\y) rectangle (4.0,\y+\h);
      \draw[c] (4.3,\y) rectangle (5.9,\y+\h);
    }
    \node at (2.0,2.72)  {A}; \node at (5.1,2.72)  {D};
    \node at (2.0,2.0)   {B}; \node at (5.1,2.0)   {E};
    \node at (2.0,0.85)  {C}; \node at (5.1,0.85)  {F};
  \end{tikzpicture}
\end{browser}
\wtnota{Stesso HTML, stessa griglia: cambia solo l'ordine di riempimento.}
\end{affianco}

Con lo stesso HTML dell'esempio precedente si ottengono A, B e C nella prima
colonna e D, E ed F nella seconda.

\subsection{Aree con nome}

È possibile assegnare un \textbf{nome} ad aree (insiemi di celle) della griglia
e collocarvi gli elementi tramite la proprietà \attr{grid-area}. È il modo più
leggibile di descrivere l'impianto di una pagina intera.

\begin{esempio}{L'impianto classico di una pagina}
\begin{affianco}[0.52]
\begin{lstlisting}[style=css, name=areas.css]
body {
  display: grid;
  grid-template-columns: 70vw 1fr;
  grid-template-rows: auto auto 1fr auto;
  grid-template-areas:
    "header  header"
    "navbar  navbar"
    "main    sidebar"
    "footer  footer";
  height: 100vh;
  margin: 0;
}

header { grid-area: header;  }
nav    { grid-area: navbar;  }
main   { grid-area: main;    }
aside  { grid-area: sidebar; }
footer { grid-area: footer;  }
\end{lstlisting}

\begin{lstlisting}[style=html, name=areas.html]
<body>
  <header>Header</header>
  <nav>Navbar</nav>
  <main>Main</main><aside>Aside</aside>
  <footer>Footer</footer>
</body>
\end{lstlisting}

\risultato
\begin{browser}[areas.html]
  \begin{tikzpicture}[font=\Lato\scriptsize,
      c/.style={line width=0.7pt, align=center}]
    \draw[c, draw=accent!60!black, fill=accent!10!white] (0,3.3) rectangle (7.4,4.0)
      node[midway, text=accent!65!black] {\code{header}};
    \draw[c, draw=orange-gradient-6, fill=orange-gradient-1!35!white] (0,2.6) rectangle (7.4,3.25)
      node[midway, text=orange-gradient-6!70!black] {\code{navbar}};
    \draw[c, draw=primary!60!black, fill=primary!10!white] (0,0.75) rectangle (5.2,2.55)
      node[midway, text=primary!65!black] {\code{main}};
    \draw[c, draw=azure-gradient-6, fill=azure-gradient-1!40!white] (5.25,0.75) rectangle (7.4,2.55)
      node[midway, text=azure-gradient-6!55!black] {\code{sidebar}};
    \draw[c, draw=accent!60!black, fill=accent!10!white] (0,0.05) rectangle (7.4,0.7)
      node[midway, text=accent!65!black] {\code{footer}};
    \node[font=\Lato\tiny, text=neutral!40!black, anchor=north] at (2.6,-0.05) {70vw};
    \node[font=\Lato\tiny, text=neutral!40!black, anchor=north] at (6.3,-0.05) {1fr};
  \end{tikzpicture}
\end{browser}
\end{affianco}

La stringa passata a \attr{grid-template-areas} è, letteralmente, un
\textbf{disegno} della pagina: ogni riga di testo è una riga della griglia, e
ripetere lo stesso nome su celle adiacenti le unisce in un'unica area.
\end{esempio}

\info{Uno strumento per esercitarsi}{
  Due giochi molto usati per imparare questi due meccanismi sono
  \urlx{https://flexboxfroggy.com/} per Flexbox e
  \urlx{https://cssgridgarden.com/} per Grid.}

\section{Media query}

Le \term{media query} permettono di applicare certi stili CSS \textbf{solo}
quando il dispositivo che sta visualizzando il contenuto ha determinate
caratteristiche. Si introducono con la parola chiave \code{@media}.

\begin{lstlisting}[style=css, name=stampa.css]
body {
  color: white;
  background: black;
}

@media print {
  body {
    color: black;
    background: transparent;
  }
}
\end{lstlisting}

Il CSS moderno supporta tre \term{tipi di output}:

\begin{center}
\renewcommand{\arraystretch}{1.4}
\begin{tabular}{@{}l@{\hspace{1.2cm}}>{\raggedright\arraybackslash}p{9.4cm}@{}}
  \textbf{Tipo} & \textbf{Destinatario} \\
  \code{print}  & materiale impaginato e documenti visualizzati in anteprima di
                  stampa. \\
  \code{screen} & dispositivi che mostrano il documento su uno schermo. \\
  \code{all}    & tutti i dispositivi di output. \\
\end{tabular}
\end{center}

\subsection{Caratteristiche del media}

Le media query possono anche verificare caratteristiche specifiche dello user
agent, del dispositivo di output o dell'ambiente. Si chiamano \term{media
feature}. La sintassi completa è:

\begin{center}
  \code{@media tipo\_media and (caratteristica: valore) and ...}
\end{center}

\begin{itemize}
  \item se il tipo è omesso, si assume \code{all};
  \item le caratteristiche vanno racchiuse fra parentesi e sono opzionali.
\end{itemize}

\begin{esempio}{Orientamento dello schermo}
\begin{affianco}[0.55]
\begin{lstlisting}[style=css, name=orientamento.css]
@media screen and (orientation: landscape) {
  h1::after { content: "Orizzontale"; }
}

@media screen and (orientation: portrait) {
  h1::after { content: "Verticale"; }
}
\end{lstlisting}
\begin{lstlisting}[style=html, name=orientamento.html]
<body><h1></h1></body>
\end{lstlisting}
\risultato
\begin{center}
  \begin{tikzpicture}[font=\Lato\scriptsize]
    \draw[draw=neutral!45!black, line width=1pt, rounded corners=3pt]
      (0,0.4) rectangle (4.4,3.1);
    \node[font=\Lora\bfseries\small] at (2.2,1.75) {Orizzontale};
    \node[anchor=north, text=neutral!35!black] at (2.2,0.3) {finestra più larga che alta};
    \draw[draw=neutral!45!black, line width=1pt, rounded corners=3pt]
      (5.1,-0.1) rectangle (7.0,3.1);
    \node[font=\Lora\bfseries\small] at (6.05,1.5) {Verticale};
    \node[anchor=north, text=neutral!35!black] at (6.05,-0.2) {più alta che larga};
  \end{tikzpicture}
\end{center}
\wtnota{L'unico contenuto della pagina è generato dal CSS, con \code{::after}.}
\end{affianco}
\end{esempio}

\begin{esempio}{Punti di rottura sulla larghezza}
\begin{lstlisting}[style=css, name=breakpoint.css]
@media (max-width: 600px)  { h1::after { content: "XSmall"; } }
@media (min-width: 600px)  { h1::after { content: "Small";  } }
@media (min-width: 768px)  { h1::after { content: "Medium"; } }
@media (min-width: 992px)  { h1::after { content: "Large";  } }
@media (min-width: 1200px) { h1::after { content: "XLarge"; } }
\end{lstlisting}

\begin{center}
  \begin{tikzpicture}[font=\Lato\scriptsize]
    \draw[-{Stealth[length=6pt]}, draw=neutral!45!black, line width=0.8pt]
      (0,0) -- (16.5,0);
    \node[anchor=north west, text=neutral!35!black] at (15.3,-0.1) {larghezza};
    \foreach \x/\v in {2.6/600px, 6.2/768px, 9.8/992px, 13.4/1200px} {
      \draw[draw=accent!70!black, line width=0.8pt] (\x,-0.18) -- (\x,0.18);
      \node[anchor=north, text=accent!70!black, font=\Lato\tiny] at (\x,-0.22) {\v};
    }
    \foreach \a/\b/\lbl in {0/2.6/XSmall, 2.6/6.2/Small, 6.2/9.8/Medium,
                            9.8/13.4/Large, 13.4/16.2/XLarge} {
      \fill[primary!12!white] (\a+0.06,0.3) rectangle (\b-0.06,0.95);
      \node[font=\Lora\small, text=primary!70!black] at ({(\a+\b)/2},0.62) {\lbl};
    }
  \end{tikzpicture}
\end{center}

Le larghezze a cui il layout cambia si chiamano \term{breakpoint} (punti di
rottura). Si noti che più regole possono essere vere contemporaneamente: a
1300px sono soddisfatte tutte tranne la prima, e prevale l'ultima in ordine di
comparsa, per la regola della cascata vista nel capitolo precedente.
\end{esempio}

\section{Responsive design}

Le media query esistono per rispondere a un problema concreto, che ha una storia
istruttiva.

\subsection{Layout a larghezza fissa}

Alla fine degli anni Novanta la maggior parte dei monitor era di 640$\times$480
pixel. Nei primi tempi del web design era quindi ragionevole progettare le
pagine per una larghezza di 640 pixel.

Poi i monitor cominciarono a crescere: 800$\times$600, poi 1024$\times$768, e
i progettisti aggiornarono di conseguenza le loro larghezze fisse.

Lavorare con una singola larghezza specifica, che sia 640, 800 o 1024 pixel, si
chiama \term{fixed-width design}. Il problema è evidente: il risultato è buono
\textbf{solo} a quella larghezza.

\begin{itemize}
  \item su schermi \textbf{più piccoli} non si riesce a vedere l'intera pagina
        senza scorrimento orizzontale;
  \item su schermi \textbf{più grandi} resta moltissimo spazio sprecato.
\end{itemize}

\subsection{Layout liquidi}

Mentre la maggioranza usava larghezze fisse, alcuni resero i propri layout più
flessibili. Non esistevano ancora né Flexbox, né Grid, né le media query: lo
fecero usando le \textbf{percentuali} come larghezza delle colonne. Questi
layout si chiamano \term{liquidi} (o \textit{fluidi}). Wikipedia fu uno dei
primi siti a usarli.

I layout liquidi funzionano bene su un intervallo ampio di larghezze, ma
peggiorano agli estremi: su uno schermo molto largo il risultato appare
\textbf{stirato}, su uno molto stretto appare \textbf{schiacciato}.

\subsection{Siti separati}

Poi arrivarono i dispositivi mobili, con schermi larghi anche solo 240 pixel,
su cui né i layout fissi né quelli liquidi funzionavano bene. Una soluzione fu
mantenere un \textbf{sito separato} per i dispositivi mobili, ospitato su un
sottodominio.

Tipicamente si ricorreva allo \term{user agent sniffing}: si esaminava l'header
\code{User-Agent} della richiesta per redirigere gli utenti mobili verso la
versione dedicata. Gli svantaggi:

\begin{itemize}
  \item lo sniffing dello user agent è \textbf{inaffidabile}: può non
        riconoscere un utente mobile, oppure redirigere per errore un utente
        desktop;
  \item ogni sito va \textbf{mantenuto separatamente};
  \item oggi la distinzione fra mobile e non mobile è sempre più sfumata.
\end{itemize}

Sarebbe stato molto meglio avere un \textbf{unico} sito che si presenta in modo
diverso a seconda delle caratteristiche del dispositivo.

\subsection{Layout adattivi}

Con l'introduzione delle media query (attorno al 2009) diventarono possibili
layout più flessibili. All'inizio, però, i progettisti erano ancora abituati
alle larghezze fisse: realizzarono quindi siti che \textbf{commutavano fra
poche larghezze fisse} usando le media query. Questo approccio si chiama
\term{adaptive design}.

\begin{affianco}[0.5]
\begin{lstlisting}[style=css, name=adattivo.css]
main  { display: inline-block; height: 200px; }
aside { display: inline-block; height: 200px; }

@media (max-width: 800px) {
  main  { width: 400px; }
  aside { width: 180px; }
}

@media (min-width: 800px) {
  main  { width: 500px; }
  aside { width: 280px; }
}
\end{lstlisting}
\risultato
\begin{browser}[finestra da 700px]
  \begin{tikzpicture}[font=\Lato\scriptsize]
    \fill[primary!14!white] (0,0) rectangle (3.4,1.3);
    \node[font=\Lora\small, text=primary!70!black] at (1.7,0.65) {main};
    \fill[azure-gradient-1!45!white] (3.5,0) rectangle (5.0,1.3);
    \node[font=\Lora\small, text=azure-gradient-6!55!black] at (4.25,0.65) {aside};
    \draw[dashed, draw=neutral!50!white] (5.1,0) rectangle (6.0,1.3);
    \node[font=\Lato\tiny, text=neutral!45!black] at (5.55,0.65) {vuoto};
  \end{tikzpicture}
\end{browser}
\begin{browser}[finestra da 1400px]
  \begin{tikzpicture}[font=\Lato\scriptsize]
    \fill[primary!14!white] (0,0) rectangle (3.0,1.3);
    \node[font=\Lora\small, text=primary!70!black] at (1.5,0.65) {main};
    \fill[azure-gradient-1!45!white] (3.1,0) rectangle (4.7,1.3);
    \node[font=\Lora\small, text=azure-gradient-6!55!black] at (3.9,0.65) {aside};
    \draw[dashed, draw=neutral!50!white] (4.8,0) rectangle (7.4,1.3);
    \node[font=\Lato\tiny, text=neutral!45!black] at (6.1,0.65) {spazio sprecato};
  \end{tikzpicture}
\end{browser}
\wtnota{Le larghezze sono fisse: oltre gli 800px il layout non cresce più con
        la finestra.}
\end{affianco}

I layout adattivi permettono di avere un solo sito che appare bene a poche
dimensioni diverse. Il risultato però non è mai davvero corretto per le
dimensioni \textbf{intermedie}: all'utente viene mostrato il layout più vicino
alla dimensione della sua finestra, e vista la varietà di dispositivi
esistenti, è probabile che la maggior parte degli utenti veda un layout meno
che ottimale.

\subsection{Layout responsive}

La sintesi arriva qui.

\begin{itemize}
  \item I layout \textbf{adattivi} sono media query + layout a larghezza fissa.
  \item I layout \textbf{responsive} sono media query + layout \textbf{liquidi}.
\end{itemize}

Il termine fu coniato da \textbf{Ethan Marcotte} in un articolo del 2010. Il
responsive design è caratterizzato da tre ingredienti:

\begin{itemize}
  \item \term{contenitori fluidi};
  \item \term{media fluidi} (immagini e video che scalano);
  \item \term{media query}.
\end{itemize}

L'impianto e le immagini di un sito responsive devono apparire bene su
\textbf{qualunque} dispositivo, non su un elenco predefinito di dimensioni.

\begin{figure}[H]
  \centering
  \begin{tikzpicture}[font=\Lato\scriptsize,
      b/.style={rounded corners=3pt, line width=0.8pt, align=center,
                minimum height=1.15cm, minimum width=2.7cm, inner sep=3pt},
      a/.style={-{Stealth[length=5pt,width=4pt]}, line width=0.9pt,
                draw=neutral!45!black}]
    \node[b, draw=neutral!55!white, fill=neutral!10!white] (f) at (0,0)
      {\textbf{Fixed-width}\\ \scriptsize una sola larghezza};
    \node[b, draw=azure-gradient-6, fill=azure-gradient-1!40!white] (l) at (3.4,0)
      {\textbf{Liquid}\\ \scriptsize percentuali};
    \node[b, draw=orange-gradient-6, fill=orange-gradient-1!40!white] (a) at (6.8,0)
      {\textbf{Adaptive}\\ \scriptsize media query + fisso};
    \node[b, draw=green-gradient-6, fill=green-gradient-1!50!white] (r) at (10.2,0)
      {\textbf{Responsive}\\ \scriptsize media query + liquido};

    \draw[a] (f) -- (l);
    \draw[a] (l) -- (a);
    \draw[a] (a) -- (r);

    \node[font=\Lato\scriptsize, text=neutral!40!black, anchor=north] at (0,-0.75)   {anni '90};
    \node[font=\Lato\scriptsize, text=neutral!40!black, anchor=north] at (3.4,-0.75) {primi 2000};
    \node[font=\Lato\scriptsize, text=neutral!40!black, anchor=north] at (6.8,-0.75) {2009};
    \node[font=\Lato\scriptsize, text=neutral!40!black, anchor=north] at (10.2,-0.75) {dal 2010};
  \end{tikzpicture}
  \caption{Breve storia dei layout web.}
  \label{fig:layout-history}
\end{figure}

\begin{esempio}{Un layout responsive}
\begin{lstlisting}[style=css, name=responsive.css]
body { display: flex; justify-content: center; flex-wrap: wrap; }

main  { display: inline-block; }
aside { display: inline-block; min-height: 100px; }
img   { width: 20%; min-width: 100px; }

@media (max-width: 600px) { main { width: 100%; } aside { width: 100%; } }
@media (min-width: 600px) { main { width:  60%; } aside { width:  40%; } }
@media (min-width: 768px) { main { width:  70%; } aside { width:  20%; } }
\end{lstlisting}

\begin{center}
  \begin{tikzpicture}[font=\Lato\scriptsize,
      win/.style={draw=neutral!45!black, line width=0.9pt, rounded corners=2pt}]
    % telefono: uno sotto l'altro
    \draw[win] (0,0) rectangle (2.4,3.4);
    \fill[primary!14!white] (0.12,1.85) rectangle (2.28,3.28);
    \node[font=\Lora\small, text=primary!70!black] at (1.2,2.56) {main};
    \fill[azure-gradient-1!45!white] (0.12,0.12) rectangle (2.28,1.73);
    \node[font=\Lora\small, text=azure-gradient-6!55!black] at (1.2,0.92) {aside};
    \node[anchor=north, text=neutral!35!black] at (1.2,-0.1) {$<$ 600px: 100\% e 100\%};

    % tablet
    \draw[win] (3.4,0) rectangle (8.2,3.4);
    \fill[primary!14!white] (3.52,0.12) rectangle (6.4,3.28);
    \node[font=\Lora\small, text=primary!70!black] at (4.96,1.7) {main};
    \fill[azure-gradient-1!45!white] (6.5,0.12) rectangle (8.08,3.28);
    \node[font=\Lora\small, text=azure-gradient-6!55!black] at (7.29,1.7) {aside};
    \node[anchor=north, text=neutral!35!black] at (5.8,-0.1) {600--768px: 60\% e 40\%};

    % desktop
    \draw[win] (9.2,0) rectangle (16.4,3.4);
    \fill[primary!14!white] (9.32,0.12) rectangle (14.4,3.28);
    \node[font=\Lora\small, text=primary!70!black] at (11.86,1.7) {main};
    \fill[azure-gradient-1!45!white] (14.5,0.12) rectangle (15.94,3.28);
    \node[font=\Lora\small, text=azure-gradient-6!55!black] at (15.22,1.7) {aside};
    \node[anchor=north, text=neutral!35!black] at (12.8,-0.1) {$>$ 768px: 70\% e 20\%};
  \end{tikzpicture}
\end{center}

La differenza rispetto all'esempio adattivo è tutta nelle unità: qui le
larghezze sono in \textbf{percentuale}, quindi fra un breakpoint e l'altro il
layout continua ad adattarsi con continuità invece di restare congelato.
\end{esempio}

\section{Il meta tag viewport}

C'è un ultimo tassello, senza il quale nulla di tutto questo funziona sui
telefoni.

I primi browser mobili dovevano gestire siti progettati per schermi molto più
larghi dei loro. Adottarono quindi una soluzione pragmatica:

\begin{enumerate}
  \item assumevano che i siti fossero progettati per uno schermo largo 980
        pixel;
  \item disegnavano la pagina in un \term{viewport virtuale} di 980 pixel;
  \item riducevano in scala il risultato per farlo entrare nella larghezza
        reale dello schermo.
\end{enumerate}

\begin{figure}[H]
  \centering
  \begin{tikzpicture}[font=\Lato\scriptsize,
      a/.style={-{Stealth[length=6pt,width=5pt]}, line width=0.9pt,
                draw=accent!70!black}]
    % viewport virtuale
    \draw[draw=primary!65!black, fill=primary!5!white, line width=0.9pt]
      (0,0) rectangle (6.2,4.0);
    \foreach \y in {3.5,3.2,2.9} \fill[primary!25!white] (0.3,\y) rectangle (5.9,\y+0.2);
    \fill[primary!18!white] (0.3,1.0) rectangle (2.9,2.6);
    \fill[primary!18!white] (3.3,1.0) rectangle (5.9,2.6);
    \node[font=\Lato\scriptsize, text=primary!65!black, anchor=north]
      at (3.1,-0.15) {viewport virtuale da 980px};

    % freccia
    \draw[a] (6.6,2.0) -- (8.0,2.0);
    \node[font=\Lato\scriptsize, text=accent!70!black, anchor=south, align=center]
      at (7.3,2.15) {riduzione\\in scala};

    % schermo reale
    \draw[draw=neutral!35!black, fill=white, line width=1.4pt, rounded corners=4pt]
      (8.5,0.35) rectangle (10.6,3.65);
    \foreach \y in {3.15,2.99,2.83} \fill[primary!25!white] (8.65,\y) rectangle (10.45,\y+0.1);
    \fill[primary!18!white] (8.65,1.6) rectangle (9.5,2.55);
    \fill[primary!18!white] (9.6,1.6) rectangle (10.45,2.55);
    \node[font=\Lato\scriptsize, text=neutral!35!black, anchor=north, align=center]
      at (9.55,0.2) {schermo reale\\480px};
  \end{tikzpicture}
  \caption{Il meccanismo del viewport virtuale dei primi browser mobili.}
  \label{fig:virtual-viewport}
\end{figure}

Questo meccanismo permetteva ai siti \textbf{non} ottimizzati per il mobile di
apparire decenti su schermi stretti. Non funziona però per i siti che
\textbf{sono} ottimizzati per il mobile: le media query impostate a 640 pixel
non scatterebbero mai, perché il browser sta disegnando in un viewport da 980.

La soluzione è il meta tag \term{viewport}, che permette di controllare questo
meccanismo. Le pagine ottimizzate per il mobile devono includere nel proprio
\tg{head}:

\begin{lstlisting}[style=html, name=viewport]
<meta name="viewport" content="width=device-width, initial-scale=1">
\end{lstlisting}

Il tag specifica due regole:

\begin{itemize}
  \item \code{width=device-width} dice al browser di assumere che la larghezza
        per cui il sito è stato progettato sia \textbf{la larghezza del
        dispositivo};
  \item \code{initial-scale=1} dice al browser di \textbf{non applicare alcuna
        riduzione} in scala.
\end{itemize}

\warning{Senza questa riga il responsive non funziona}{
  È l'errore più frustrante da diagnosticare: si scrive un foglio di stile
  responsive impeccabile, lo si prova ridimensionando la finestra del computer e
  funziona benissimo; poi lo si apre sul telefono e appare tutto minuscolo, come
  una pagina desktop rimpicciolita. La causa è quasi sempre la mancanza del meta
  tag viewport nell'\tg{head}.}

\section{Riferimenti}

\begin{itemize}
  \item \textbf{Learn CSS} --- web.dev. Sezioni 2, da 8 a 11. \\
        \urlx{https://web.dev/learn/css/}
  \item \textbf{Learn Responsive Design} --- web.dev. Sezioni da 1 a 3 e 15. \\
        \urlx{https://web.dev/learn/design/}
  \item \textbf{Flexbox Froggy} --- gioco per imparare Flexbox. \\
        \urlx{https://flexboxfroggy.com/}
  \item \textbf{CSS Grid Garden} --- gioco per imparare Grid. \\
        \urlx{https://cssgridgarden.com/}
\end{itemize}
